\documentclass[12pt]{article}
\usepackage[T1]{fontenc} 
\usepackage[utf8x]{inputenc} 
\usepackage{lmodern}
\usepackage [frenchb]{babel}
% Pour pouvoir utiliser 
\usepackage{ucs}

\usepackage{textcomp}
\usepackage{graphicx}
\usepackage{keystroke}
\usepackage{amssymb}
\usepackage{amsmath}
%\renewcommand{\thesection}{\arabic{section}} % numérotation des sections	
\usepackage[cc]{titlepic} %rajouter le logo dans la page de garde
\usepackage{url} % Pour avoir de belles url
\usepackage[a4paper]{geometry}
\usepackage[linktocpage]{hyperref}
% Pour pouvoir faire plusieurs colonnes
\usepackage {multicol}
%Pour faire plusieurs lignes
\usepackage{multirow}
\usepackage{slashbox}
% Pour mettre du code source
\usepackage {listings}
% Pour pouvoir passer en paysage
\usepackage{lscape}	

% Pour pouvoir faire plusieurs colonnes
\usepackage {multicol}

% POur crééer un index
\usepackage{makeidx}

%insertion de pdf
\usepackage{pdfpages}

%Utilisation d'acronyme
\usepackage{acronym}

\usepackage{fancyhdr}

\pagestyle{fancy}
\usepackage{lastpage} %numérotation type k/n
\renewcommand{\footrulewidth}{0.5pt} % trait horizontal
%\renewcommand{\headrulewidth}{1pt} % suppresion du trait horizontal dans l'entête
\lfoot{Analyse of a research article}
\cfoot{\thepage/2 } % numérotation des pages
\rfoot{A.\textsc{M} \& M.\textsc{O}}
%\fancyhead[L]{\hspace{-1.5cm}\includegraphics[scale=0.105]{img/logouniv.eps}} %logo entête !
\fancyhead{\empty} %logo entête !


\fancypagestyle{plain}{% 1ères pages des chapitres:
   \fancyhead{} % supprime l’entete...
   \renewcommand{\headrulewidth}{0pt} % ...et le filet
}

\renewcommand{\thesection}{\arabic{section}}





\hypersetup{
  backref=true,
  %permet d'ajouter des liens dans...
  pagebackref=true,%...les bibliographies
  hyperindex=true, %ajoute des liens dans les index.
  colorlinks=true, %colorise les liens
  breaklinks=true, %permet le retour à la ligne dans les liens trop longs
  urlcolor= blue, %couleur des hyperliens
  citecolor=	cyan,
  bookmarks=true, %créé des signets pour Acrobat
  bookmarksopen=true,
  %si les signets Acrobat sont créés,
  %les afficher complètement.
  pdftitle={ScProjectArticle}, %informations apparaissant dans
  pdfauthor={AM and MO},
  %dans les informations du document
  pdfsubject={Doc}
  %sous Acrobat.
}

%\makeindex
%%%% debut macro pour enlever le nom chapitre %%%%
\makeatletter
\def\@makechapterhead#1{%
  \vspace*{50\p@}%
          {\parindent \z@ \raggedright \normalfont
            \interlinepenalty\@M
            \ifnum \c@secnumdepth >\m@ne
            \Huge\bfseries \thechapter\quad
            \fi
            \Huge \bfseries #1\par\nobreak
            \vskip 40\p@
}}

\def\@makeschapterhead#1{%
  \vspace*{50\p@}%
          {\parindent \z@ \raggedright
            \normalfont
            \interlinepenalty\@M
            \Huge \bfseries  #1\par\nobreak
            \vskip 40\p@
}}
\makeatother
%%%% fin macro %%%%
%Couverture 
\widowpenalty=10000
\clubpenalty=10000


\title{ {\huge \textsc{Senior capstone project}} \\  Master ALMA  \\ \textbf{Analyse of a research article} \\ \emph{Characterizing Software Architecture Changes : A Systematic Review}\\{\small Supervisor :  M.\textsc{Oussalah}}}
\author{A.\textsc{Marguerite} \& M.\textsc{Ouairy}}
\newcommand{\cf}{\emph{cf}.}


\begin{document}

\maketitle
\renewcommand{\labelitemi}{$\bullet$} 



\section{Introduction}\label{sec:intro}
This article was written by Byron J. \textsc{Williams} who is an assistant professor of the Department of Computer Science at the \href{http://www.msstate.edu/}{Mississippi State University} and Jeffrey C. \textsc{Carver} assistant professor of the Department of Computer Science at the \href{http://www.ua.edu/}{University of Alabama}. The research focus area of this article are about changes and theirs impact in architecture of information system. In the first part of this article, the authors put forward a systematic literature review, that identifies and characterizes software changes (\cf{} section \ref{sec:pblm}). Then, in a second part, they introduce the \emph{Software Architecture Change Characterization Scheme} (SACCS). We will go into this contribution in greater depth in the section \ref{sec:res}. 


%% Cet article a été rédigé par M.\textsc{Butler} chercheur à l'\href{http://www.southampton.ac.uk/}{Université de Southampton} (UK) et K.\textsc{Damchoom} actuellement en Postdoc  à l'université de \href{http://www.southampton.ac.uk/}{Université de Southampton} (UK). Le travail proposé par ces deux chercheurs porte autour des méthodes de spécifications et de raisonnements formels. Plus particulièrement, ils utilisent la méthode \emph{Event-B}. Il s'agit d'une évolution de la méthode B introduisant la notion d'événements. Le remplacement des \og opérations \fg{} de la méthode B par des actions dotée1s de gardes donne la possibilité d'étudier des systèmes orientés plus \emph{hardware}. C'est la démarche choisie par les auteurs  dont la problématique sera développée dans la section \ref{sec:pblm}. 

\section{Research area}\label{sec:pblm}
Nowadays, developers know an increasing demands on software. These demands involve new software as well the upgrades of existing ones. Such operations lead developers to face up with complexities and the risks of degrading the software architecture. A systematic review is a repeatable and formalized process in which researcher systematically search a body of literature to document the state of knowledge on a particular subject. A such report, give to researchers who use it a better hindsight than an ad-hoc review.% After this contribution done, \textsc{Williams} and \textsc{Carver} 

%% M.\textsc{Butler} et K.\textsc{Damchoom} proposent dans ce papier d'utiliser les méthodes décrites dans la section \ref{sec:intro} avec un exemple d'application concrète. Il s'agit de l'utilisation d'une mémoire flash, et de la gestion d'erreurs de ses opérations. Par exemple comment éviter la perte de données lors d'une coupure de courant, ou suite à un retrait manuel au cours d'une opération ? Suite à des éventuelles erreurs, quels mécanismes mettre en place pour garantir la consistance des données ? Les auteurs proposent d'étudier ce problème à partir d'une architecture flash (\href{http://www.onfi.org}{ONFI}) en mettant en oeuvre les outils proposés par la méthode Event-B avec l'aide de l'IDE (Rodin) associé.


\section{Main results}\label{sec:res}
\textsc{Williams} and \textsc{Carver} followed the classical steps to write this systematic review. After the identification of the needs (\cf{} section \ref{sec:pblm}), they established several very restrictive questions to extract only relevant data. This step was very important because of the size of the database they used. Indeed, by using several databases like \textit{IEEE Electronic Library} or \textit{Google Scholar}, they were looked for any journal articles, conference proceeding or book chapter that dealt with their selected criterion. Thus, every question selected by the authors is justified by a solid motivation. Thanks to this process, \textsc{Williams} and \textsc{Carver} filtered 2752 papers among thousands. Then they started a manual selection by reading the title (in a first move) then the abstract (in a second move) of each paper. To conclude this data extraction the authors used several inclusion and exclusion criterion. Contrary to the first restrictive questions which only concerned the content (link with the research topic), these criterion concerned also the form of papers. With the 130 papers left, the authors could begin the summary and the analysis to interpret their results and then answer to first questions.

\paragraph{}
The \textit{SACCS} was developed by conducting the systematic literature described earlier. The main goal of this scheme is to help a developer when a change request appears. It is composed by a set of characteristics and their attribute associated. To use the \textit{SACCS}, a developer has to pick a high-level characteristic according to his change request. Then, he could get more information by progressing into the scheme. For each value of an attribute, the authors give a description and the list of the papers which deal with this topic.

\section{Discussion an re-use}
In conclusion, the writing of a systematic review of software change leads the authors to the production of the \textit{SACCS}. The main difference to others \textit{change classification schemes} is that developer is assisted in making decision thanks to many important characteristics including developer's experience. Therefore, the knowledge of the change will help the developer to not degrading the quality of the software. It can be noted that the writing of a systematic review produces on one hand a very useful tool for the researcher's community, and on the other hand a way to find and offer a new contribution.

\end{document}

%%  LocalWords:  Butler Southampton UK Damchoom Postdoc Event-B ONFI
%%  LocalWords:  l'IDE Fourth refinement skip end fail compt write
%%  LocalWords:  begin step multi temp SACCS hoc identificate IEEE
% LocalWords:  Google english informations
